\section{Sequential Regime Complexity}\label{sec:sequential}

We now move to sequential settings with transitions and observations (POMDP-style). Sufficiency becomes a policy-level invariance condition.

\subsection{Sequential Decision Problems}

\begin{definition}[Sequential Decision Problem]\label{def:sequential-decision-problem}
A sequential decision problem is a tuple
\[
\mathcal{D}_{\mathrm{seq}}=(A,X_1,\ldots,X_n,U,T,O),
\]
with state space $S=X_1\times\cdots\times X_n$, transition kernel $T:A\times S\to\Delta(S)$, and observation model $O:S\to\Delta(\Omega)$.
\end{definition}

\begin{definition}[Sequential Sufficiency]\label{def:sequential-sufficient}
Coordinate set $I$ is \emph{sequentially sufficient} if the optimal policy using full observation history equals the optimal policy obtained when histories are restricted to $I$-coordinates.
\end{definition}

\begin{problem}[SEQUENTIAL-SUFFICIENCY-CHECK]\label{prob:sequential-sufficiency}
\textbf{Input:} $\mathcal{D}_{\mathrm{seq}}$ and $I\subseteq\{1,\ldots,n\}$. \\
\textbf{Question:} Is $I$ sequentially sufficient?
\end{problem}

\subsection{PSPACE-Completeness}

\begin{theorem}[SEQUENTIAL-SUFFICIENCY-CHECK is PSPACE-complete]\label{thm:sequential-pspace}
SEQUENTIAL-SUFFICIENCY-CHECK is PSPACE-complete.
\end{theorem}

\begin{proof}
Membership: finite-horizon policy comparison for POMDP-style models is decidable in polynomial space by dynamic programming over belief/state summaries.

Hardness: reduce TQBF. Alternating quantifiers are encoded as alternating controlled/adversarial transitions; terminal reward evaluates the quantified formula. Full-history sufficiency holds exactly when the quantified instance is true.
\end{proof}

\subsection{Tractable Subcases}

\begin{proposition}[Sequential Tractable Classes]\label{prop:sequential-tractable}
SEQUENTIAL-SUFFICIENCY-CHECK is polynomial-time solvable under each of the following restrictions:
\begin{enumerate}
\item full observability (MDP case),
\item bounded horizon,
\item tree-structured transition systems,
\item deterministic transitions.
\end{enumerate}
\end{proposition}

\begin{proof}
Each restriction reduces policy search to dynamic programming over a polynomially bounded representation.
\end{proof}

\subsection{Bridge from Stochastic}

\begin{proposition}[Stochastic Sufficiency Does Not Imply Sequential Sufficiency]\label{prop:stochastic-not-sequential}
There exist instances where $I$ is sufficient in the stochastic one-shot sense but fails to be sufficient once temporal dynamics are introduced.
\end{proposition}

\begin{proof}
Construct a process in which the same one-step expected optimizer is preserved under $I$, but hidden-state memory affects future information states and therefore optimal policies. The temporal dependence breaks one-shot sufficiency transfer.
\end{proof}
